% Part: computability
% Chapter: computability-theory
% Section: fixed-point-thm

\documentclass[../../../include/open-logic-section]{subfiles}

\begin{document}

\olfileid[tr]{cmp}{thy}{fix}
\olsection{Sabit Nokta Teoremi}

Durma problemini yeniden ele alalım. Geçici bir gösterim olarak
$\tuple{x,y}$ yerine $\gn{\cfind{x}(y)}$ yazalım ve bunu
$\cfind{x}(y)$ değerinin bir ``adı'' olarak düşünelim. Bu gösterimle durma
probleminin karar verilemez olduğuna ilişkin ispatlarımızdan birini yeniden
ifade edebiliriz.

Soru: Her $x$ ve $y$ için
\[
  h(\gn{\cfind{x}(y)})=\begin{cases}
    1 & \text{$\cfind{x}(y)\fdefined$ ise}\\
    0 & \text{aksi durumda}
  \end{cases}
\]
özelliğini taşıyan hesaplanabilir bir $h$ fonksiyonu var mıdır?

Yanıt: Hayır. Olsaydı
\[
  g(x)\simeq\begin{cases}
    0 & \text{$h(\gn{\cfind{x}(x)})=0$ ise}\\
    \text{tanımsız} & \text{aksi durumda}
  \end{cases}
\]
kısmi fonksiyonu hesaplanabilir olur ve dolayısıyla bir $e$ indisine sahip
olurdu. Fakat o zaman
\[
  \cfind{e}(e)\simeq\begin{cases}
    0 & \text{$h(\gn{\cfind{e}(e)})=0$ ise}\\
    \text{tanımsız} & \text{aksi durumda}
  \end{cases}
\]
olur ve $\cfind{e}(e)$ ancak ve ancak tanımlı değilse tanımlı olurdu; bu bir
çelişkidir.

Şimdi $\cfind{e}$ içeren denkleme bakın. Bir anlamda burada kendine gönderme
vardır: $\cfind{e}(e)$ değerinin belirli bir yolla
$\gn{\cfind{e}(e)}$ ifadesine bağlı olmasını sağladık. Sabit nokta teoremi,
bunu yalnızca çelişki ispatlamak amacıyla değil, genel olarak
\emph{yapabileceğimizi} söyler.

\olref{lem:fixed-equiv}, sabit nokta teoremini ifade etmenin iki eşdeğer
yolunu verir. Mantıksal bakımdan ifadelerin eşdeğerliği ikisinin de doğru
olmasından çıkar; fakat asıl kastımız her birinin ötekinden doğrudan çıkması
ve böylece aynı teoremin alternatif ifadeleri olarak alınabilmeleridir.

\begin{lem}
\ollabel{lem:fixed-equiv}
Aşağıdaki ifadeler eşdeğerdir:
\begin{enumerate}
\item Her kısmi hesaplanabilir $g(x,y)$ fonksiyonu için öyle bir $e$ indisi
  vardır ki her $y$ için
  \[
    \cfind{e}(y)\simeq g(e,y)
  \]
  olur.
\item Her hesaplanabilir $f(x)$ fonksiyonu için öyle bir $e$ indisi vardır
  ki her $y$ için
  \[
    \cfind{e}(y)\simeq\cfind{f(e)}(y)
  \]
  olur.
\end{enumerate}
\end{lem}

\begin{proof}
$(1)\Rightarrow(2)$: $f$ verilsin ve
$g(x,y)\simeq\fn{Un}(f(x),y)$ olarak tanımlansın. (1)'i kullanarak her $y$
için
\begin{align*}
  \cfind{e}(y) & = \fn{Un}(f(e),y)\\
  & = \cfind{f(e)}(y)
\end{align*}
olacak bir $e$ indisi elde edin.

$(2)\Rightarrow(1)$: $g$ verilsin. $s$-$m$-$n$ teoremini kullanarak her
$x$ ve $y$ için $\cfind{f(x)}(y)\simeq g(x,y)$ olacak bir $f$ elde edin.
(2)'yi kullanarak
\begin{align*}
  \cfind{e}(y) & = \cfind{f(e)}(y)\\
  & = g(e,y)
\end{align*}
olacak bir $e$ indisi elde edin. Bu, ispatı tamamlar.
\end{proof}

\begin{explain}
(1) ifadesinin doğru olduğunu, dolayısıyla (2)'nin de doğru olduğunu
göstermeden önce ne kadar şaşırtıcı olduğunu düşünün. $e$'yi bir bilgisayar
programı olarak düşünürsek (1), herhangi bir kısmi hesaplanabilir $g(x,y)$
verildiğinde $g_e(y)\simeq g(e,y)$ fonksiyonunu hesaplayan bir $e$ programı
bulabileceğinizi söyler. Başka bir deyişle kendisine gönderme yapan bir
fonksiyonu hesaplayan bilgisayar programı bulabilirsiniz.
\end{explain}

\begin{thm}
\olref{lem:fixed-equiv} içindeki iki ifade doğrudur. Özellikle her kısmi
hesaplanabilir $g(x,y)$ fonksiyonu için öyle bir $e$ indisi vardır ki her
$y$ için
\[
  \cfind{e}(y)\simeq g(e,y)
\]
olur.
\end{thm}

\begin{proof}
Bileşenler yukarıdaki durma problemi tartışmasında zaten örtük olarak
bulunur. $\fn{diag}(x)$, her $x$ için
$f_x(y)\simeq\cfind{x}(x,y)$ fonksiyonunun bir indisini döndüren
hesaplanabilir bir fonksiyon olsun; yani
\[
  \cfind{\fn{diag}(x)}(y)\simeq\cfind{x}(x,y).
\]
$\fn{diag}$ fonksiyonunu, iki yerli bir fonksiyon programını ilk argümanı
olarak özgün programı sabitleyerek elde edilen tek yerli fonksiyonun
programına dönüştüren bir fonksiyon olarak düşünün. $\fn{diag}$ biçimsel
olarak şöyle tanımlanabilir: Önce
\[
  s(x,y)\simeq\fn{Un}^2(x,x,y)
\]
ile $s$'yi tanımlayın; burada $\fn{Un}^2$, iki yerli kısmi hesaplanabilir
fonksiyonlar için evrensel olan üç yerli bir fonksiyondur. Ardından
$s$-$m$-$n$ teoremi uyarınca
\[
  \cfind{\fn{diag}(x)}(y)\simeq s(x,y)
\]
koşulunu sağlayan ilkel özyinelemeli bir $\fn{diag}$ fonksiyonu bulabiliriz.

Şimdi $l$ fonksiyonunu
\[
  l(x,y)\simeq g(\fn{diag}(x),y)
\]
ile tanımlayalım ve $\gn{l}$, $l$ için bir indis olsun. Son olarak
$e=\fn{diag}(\gn{l})$ diyelim. O zaman her $y$ için
\begin{align*}
  \cfind{e}(y) & \simeq\cfind{\fn{diag}(\gn{l})}(y)\\
  & \simeq\cfind{\gn{l}}(\gn{l},y)\\
  & \simeq l(\gn{l},y)\\
  & \simeq g(\fn{diag}(\gn{l}),y)\\
  & \simeq g(e,y)
\end{align*}
olur; istenen buydu.
\end{proof}

\begin{explain}
Burada ne oluyor? Kendisini çıktı olarak veren bir bilgisayar programı yazma
görevinin size verildiğini düşünün. Ayrıca dizge fonksiyonlarından oluşan
zengin ve tuhaf bir kitaplığa sahip bir programlama diliyle çalıştığınızı
varsayın. Özellikle programlama dilinizde şu şekilde çalışan bir
$\fn{diag}$ fonksiyonu bulunsun: $s$ girdi dizgesi verildiğinde $\fn{diag}$,
$s$ içinde geçen her `x' simgesini bulup özgün dizgenin tırnak içine alınmış
bir sürümüyle değiştirir. Örneğin girdi olarak
\begin{quote}
\begin{verbatim}
hello x world
\end{verbatim}
\end{quote}
dizgesi verildiğinde çıktı olarak
\begin{quote}
\begin{verbatim}
hello 'hello x world' world
\end{verbatim}
\end{quote}
döndürür. Bu durumda istenen programı yazmak kolaydır; şu kodun işe
yaradığını denetleyebilirsiniz:
\begin{quote}
\begin{verbatim}
print(diag('print(diag(x))'))
\end{verbatim}
\end{quote}
C++ ve Java gibi daha yaygın programlama dillerinde aynı düşünce, daha
ayrıntılı bir uygulamayla yine çalışır.

Sabit nokta teoreminin ispatına yalnızca birkaç adım uzaktayız. Yazdırma
fonksiyonunun bir $x$ dizgesi ile başka bir sayısal $y$ argümanı alan ve
$x$ dizgesini $y$ kez yazdıran $\fn{print}(x,y)$ sürümünü düşünelim. O zaman
\begin{quote}
\begin{verbatim}
getinput(y);
print(diag('getinput(y); print(diag(x), y)'), y)
\end{verbatim}
\end{quote}
``programı'', $y$ girdisinde kendisini $y$ kez yazdırır.
$\fn{getinput}$--$\fn{print}$--$\fn{diag}$ iskeletini keyfî bir $g(x,y)$
fonksiyonuyla değiştirdiğimizde
\begin{quote}
\begin{verbatim}
g(diag('g(diag(x), y)'), y)
\end{verbatim}
\end{quote}
elde edilir. Bu, $y$ girdisinde $g$'yi programın kendisi ve $y$ üzerinde
çalıştıran bir programdır. ``Tırnak içine alma''yı ``bir indis kullanma''
olarak düşünürsek yukarıdaki ispatı elde ederiz.

Şimdilik ispatı biçimsel bir hile ya da kara büyü gibi düşünebilirsiniz.
Fakat yukarıda verilen savın ayrıntılarını yeniden kurabilmelisiniz.
Eksiklik teoremlerini ve ilişkili ``sabit nokta teoremi''ni ispatladığımızda
neden çalıştığını anlamanın başka yollarını tartışacağız.
\end{explain}

\begin{tagblock}{lambda}
\begin{digress}
Aynı düşünce bir ``sabit nokta'' birleşimleyicisi elde etmek için de
kullanılabilir. Bir lambda terimi $g$ verilsin ve $k$'nin $gk$ ile
$\beta$-eşdeğer olduğu bir $k$ terimi istensin. Gösterim uzlaşımlarımızla
\[
  \fn{diag}(x)=xx
\]
ve
\[
  l(x)=g(\fn{diag}(x))
\]
terimlerini tanımlayalım; başka bir deyişle $l$, $\lambd[x][g(xx)]$
terimidir. $k$, $ll$ terimi olsun. O zaman
\begin{align*}
  k & = (\lambd[x][g(xx)])(\lambd[x][g(xx)])\\
  & \red g((\lambd[x][g(xx)])(\lambd[x][g(xx)]))\\
  & = gk
\end{align*}
olur. Eğer
\[
  Y=\lambd[g][((\lambd[x][g(xx)])(\lambd[x][g(xx)]))]
\]
alınırsa $Yg$ ile $g(Yg)$ ortak bir terime indirgenir; dolayısıyla
$Yg\equiv_\beta g(Yg)$ olur. Buna ``Curry birleşimleyicisi'' denir. Bunun
yerine
\[
  Y=(\lambd[xg][g(xxg)])(\lambd[xg][g(xxg)])
\]
alınırsa gerçekten $Yg$, $g(Yg)$'ye indirgenir; bu daha güçlü bir ifadedir.
$Y$'nin bu son sürümüne ``Turing birleşimleyicisi'' denir.
\end{digress}
\end{tagblock}

\end{document}

